\section{Function reconstruction for compressed actual return operators}
\label{sec:compressed-return-resolvents}
A distributional spectral vector does not automatically belong to a
bounded-variation class. We give a setting where the reconstruction
and the defect comparison can both be proved: finite-rank orthogonal
compressions of the actual return operator. The multiplier uses the
full unbounded return record. Only the function space is compressed.
We keep the mesh dependence explicit and prove a separate finite-time
comparison with the uncompressed characteristic function.

\subsection{An explicit regular function class}
For a sufficiently small mesh size $h>0$, subdivide the five disjoint
coordinate rectangles making up $Y_R^*$ into rectangles with side
lengths between $h/2$ and $h$. Equal subdivisions of each side have
this property. Angular coordinates are taken in a chart through each
rectangle; no artificial small cell is created at the angular seam.
Denote the partition by $\mathcal P_{R,h}$, its piecewise constant
space by $V_{R,h}\subset L^2(\nu_R^*)$, and its conditional expectation
by $P_{R,h}$. The side lengths of the original rectangles are fixed;
only their locations move with $R$.

\begin{lemma}[Uniform reconstruction on the section partition]
\label{lem:section-grid-reconstruction}
Uniformly in $R$ and $h$,
\begin{equation}\label{eq:grid-reconstruction-budget}
 \dim V_{R,h}\le C h^{-2},\qquad
 \|f\|_\infty+\|f^0\|_{\mathrm{BV}}
                         \le C h^{-1}\|f\|_{L^2(\nu_R^*)}
                    \quad(f\in V_{R,h}).
\end{equation}
The map $P_{R,h}$ is an orthogonal projection, preserves constants,
and contracts both $L^2$ and $L^\infty$. For every bounded section
function $v$ with zero extension of bounded variation,
\begin{align}
 \|(I-P_{R,h})v\|_{L^1(\nu_R^*)}
                 &\le Ch\|v^0\|_{\mathrm{BV}},
                       \label{eq:grid-L1-approximation}\\
 \|(I-P_{R,h})v\|_{L^2(\nu_R^*)}^2
                 &\le Ch\|v\|_\infty\|v^0\|_{\mathrm{BV}}.
                       \label{eq:grid-L2-approximation}
\end{align}
The zero extension in \eqref{eq:grid-reconstruction-budget} includes
all jumps at the outer section boundary.
\end{lemma}
\begin{proof}
The normalized collision density on the section is the constant
$(4\pi c_*)^{-1}$ in $(\alpha,p)$ coordinates. Each cell has mass
between $c h^2$ and $C h^2$, and there are at most $C h^{-2}$ cells.
If $f$ has value $a_B$ on cell $B$ and norm one, then
$\sum_B|a_B|^2\le C h^{-2}$ and $\max_B|a_B|\le C h^{-1}$.
Every cell side has length at most $h$. Summing internal jumps and
outer traces bounds the first variation by
\[
 C h\sum_B|a_B|
 \le C h(\#\mathcal P_{R,h})^{1/2}
                       \left(\sum_B|a_B|^2\right)^{1/2}
 \le C h^{-1}.
\]
Real and imaginary parts give the stated complex BV norm. The $L^1$
part of the norm is uniformly bounded, proving
\eqref{eq:grid-reconstruction-budget}.

For a smooth $v$, compare its value to its cell average by integrating
$|v(x)-v(y)|$ over pairs in that cell. Change the two coordinates one
at a time and integrate the fundamental theorem of calculus. Since
the aspect ratios are bounded, the result on a cell is at most $Ch$
times the integral of its first derivatives. BV approximation, or
the same one-dimensional slicing argument for distributional
variation, gives \eqref{eq:grid-L1-approximation}. Summation over
cells and the fixed normalization leave the constant uniform.
Finally $|(I-P)v|\le2\|v\|_\infty$; multiplication of the $L^1$
bound gives \eqref{eq:grid-L2-approximation}. The projection and
contraction properties follow directly from conditional expectation.
\end{proof}

Define the actual twisted Koopman operator and its compression by
\begin{equation}\label{eq:actual-compressed-operator}
 \mathcal U_{R,z}f=e^{-iz\cdot g_R}f\circ F_R^*,\qquad
 A_{R,h,z}=P_{R,h}\mathcal U_{R,z}|_{V_{R,h}}.
\end{equation}
The first operator is unitary on $L^2(\nu_R^*)$: $F_R^*$ is invertible
and measure preserving, and the real-frequency multiplier has modulus
one. Thus $\|A_{R,h,z}\|\le1$. In the orthonormal cell basis
$e_B=\one_B/\sqrt{\nu_R^*(B)}$, its exact matrix entries are
\begin{equation}\label{eq:actual-compressed-matrix}
 (A_{R,h,z})_{BC}=
 \frac{\displaystyle\int_{B\cap(F_R^*)^{-1}C}
                         e^{-iz\cdot g_R}\dd\nu_R^*}
      {\sqrt{\nu_R^*(B)\nu_R^*(C)}}.
\end{equation}
These integrals retain every return length; they are not truncated
flight records or matrices from an independent stochastic model.

\begin{lemma}[Compression residual versus physical defect]
\label{lem:compression-defect-comparison}
Let $U$ be an isometry of a Hilbert space, $P$ an orthogonal projection,
and $A=PU|_{\operatorname{ran}P}$. If $\|f\|=1$, $Pf=f$,
$|\lambda|=1$, and $\varepsilon=\|(A-\lambda)f\|\le1$, then
\begin{equation}\label{eq:compression-defect-comparison}
 \|(U-\lambda)f\|^2
   =\|(A-\lambda)f\|^2+1-\|Af\|^2
   \le3\varepsilon.
\end{equation}
In particular, a physical lower bound
$\|(U-\lambda)f\|\ge d$ for all normalized $f\in\operatorname{ran}P$,
with $0<d\le1$, gives
\begin{equation}\label{eq:compression-lower-singular}
 \|(A-\lambda)f\|\ge(d^2/3)\|f\|.
\end{equation}
\end{lemma}
\begin{proof}
The orthogonal decomposition of $(U-\lambda)f$ has components
$(A-\lambda)f$ and $(I-P)Uf$. The latter has squared norm
$1-\|Af\|^2$. Also $\|Af\|\ge1-\varepsilon$, so
$1-\|Af\|^2\le2\varepsilon-\varepsilon^2$ when
$\varepsilon\le1$. This proves the displayed, slightly weaker,
constant $3$. If $\varepsilon>1$, the lower bound in
\eqref{eq:compression-lower-singular} is immediate. Homogeneity
finishes the proof. The argument controls arbitrary approximate
vectors, not only exact eigenvectors, and makes no normality assumption
on the compression.
\end{proof}

\begin{theorem}[Quantitative peripheral resolvents of the compression]
\label{thm:compressed-peripheral-resolvent}
Choose $H_h=C_0h^{-1}\ge2$, with $C_0$ large enough for
Lemma~\ref{lem:section-grid-reconstruction}. Set
\[
 \rho=(|z|^2+\xi^2)^{1/2},\quad K_h=K(H_h,\rho),\quad
 \delta_h=b_0\rho^4/K_h^2.
\]
There are uniform $a,\rho_0,b_0>0$ such that, whenever
\begin{equation}\label{eq:compressed-peripheral-domain}
 \xi\in[-\pi,\pi],\quad0<\rho\le\rho_0,\quad \rho^2K_h\le a,
\end{equation}
the finite matrix $e^{i\xi}-A_{R,h,z}$ is invertible and
\begin{equation}\label{eq:compressed-resolvent}
 \|(e^{i\xi}-A_{R,h,z})^{-1}\|_{L^2\to L^2}
                                     \le\delta_h^{-1}.
\end{equation}
For $u\ge1$ the same bound holds at $ue^{i\xi}$. It holds with
constant $2\delta_h^{-1}$ for $1-\delta_h/2\le u\le1$.
An eigenvalue $\lambda=ue^{i\xi}$ satisfying the angular and physical
conditions in \eqref{eq:compressed-peripheral-domain} obeys
\begin{equation}\label{eq:compressed-eigenvalue-gap}
 1-u\ge\delta_h.
\end{equation}
The adjoint compression has the corresponding resolvent bound at
$e^{-i\xi}$.
\end{theorem}
\begin{proof}
Every normalized vector of $V_{R,h}$ satisfies the actual section
function budget $H_h$. Also
\[
 \|\mathcal U_{R,z}f-e^{i\xi}f\|_2
                         =\|\mathcal E_{R,z,\xi}f\|_2.
\]
Theorem~\ref{thm:joint-return-spectral-defect} therefore supplies
$d=b\rho^2/K_h$, reduced if necessary so that $d\le1$.
Apply Lemma~\ref{lem:compression-defect-comparison} and choose
$b_0=b^2/3$. Finite dimensionality turns the lower singular-value
bound into invertibility and \eqref{eq:compressed-resolvent}.

For a normalized $f$ and $u\ge1$,
\begin{align*}
 &\|(ue^{i\xi}-A)f\|^2-\|(e^{i\xi}-A)f\|^2\\
 &\qquad=(u-1)\bigl(u+1-2\operatorname{Re}
                         (e^{-i\xi}\langle f,Af\rangle)\bigr)\ge0,
\end{align*}
where $\langle f,g\rangle=\int\overline f g\dd\nu_R^*$ and
$|\langle f,Af\rangle|\le1$.
For the inner radial interval use the reverse triangle inequality;
the lower singular value is at least $\delta_h-(1-u)\ge\delta_h/2$.
If $Af=ue^{i\xi}f$, the singular-value estimate at $e^{i\xi}$ gives
$1-u\ge\delta_h$. Finally take the adjoint of the inverse to obtain
the last assertion. All statements are resolvent bounds for the
specified compression, not a spectral assertion about $\mathcal U_{R,z}$
on the full $L^2$ space.
\end{proof}

\subsection{Finite-time consistency with the unmodified record}
To relate this operator to the raw problem, we record an error bound
which is independent of a spectral convergence assertion. Use
$U_{j,R}=J_{j,R}-j\bar G_R$ and $N_{0,R}=0$.

\begin{lemma}[Variation after a collision-count cutoff]
\label{lem:truncated-return-phase-variation}
For integers $1\le j\le n\le L$ the zero extension of
\[
 v_{j,L}(x)=\one_{\{N_{j,R}(x)\le L\}}
                                  e^{-iz\cdot U_{j,R}(x)},\quad x\in Y_R^*,
\]
satisfies, uniformly in $R$,
\begin{equation}\label{eq:truncated-return-phase-variation}
 \|v_{j,L}\|_\infty\le1,\qquad
 \|v_{j,L}^0\|_{\mathrm{BV}}
                     \le(1+|z|)\exp\{C L\log(2+L)\}.
\end{equation}
\end{lemma}
\begin{proof}
Split according to the collision time $m=N_{j,R}$, $j\le m\le L$.
Starting in $Y_R^*$, this event says that
$\eta_R\circ T_R^m=1$ and exactly $j-1$ of the indicators
$\eta_R\circ T_R^i$, $1\le i<m$, equal one.
Write it as a disjoint union of the at most $2^{m-1}$ binary
membership patterns, multiplied by the initial section indicator.
Each pattern is a product of at most $m+1$ factors from
$\eta_R\circ T_R^i$ and $1-\eta_R\circ T_R^i$.
On this event exact compensation gives $U_{j,R}=S_{m,R}h_R$.

Lemma~\ref{lem:finite-record-variation} and the BV product inequality
bound the variation of each pattern by $C(m+1)B_L$. The phase
$e^{-iz\cdot S_mh_R}$ has supremum one and variation at most
$C(1+|z|m)B_L$, by the BV chain inequality. Sum their product bounds
over $m\le L$ and all patterns. The factors $2^L$ and the polynomial
powers of $L$ are absorbed by increasing the constant in $B_L$.
All section and itinerary jumps, including the initial boundary, have
been counted. No variation bound for an inverse-coarea density is
used.
\end{proof}

\begin{proposition}[An exact telescoping comparison of powers]
\label{prop:compressed-finite-time-comparison}
There are $C,C_1,a,c>0$ such that for $n\ge2$, $L\ge n$,
\begin{align}
 &\left|\int_{Y_R^*}e^{-iz\cdot U_{n,R}}\dd\nu_R^*
             -\langle1,A_{R,h,z}^n1\rangle\right|\notag\\
 &\quad\le Cn\left[
       e^{(an-cL)/2}
       +\sqrt{h(1+|z|)\exp\{C_1L\log(2+L)\}}
                    \right].
                 \label{eq:compressed-finite-time-error}
\end{align}
In particular the matrix pairing approximates the actual characteristic
function at every fixed return time when the cutoff and mesh are
chosen so that the right side tends to zero.
\end{proposition}
\begin{proof}
Abbreviate $\mathcal U_{R,z}$, $P_{R,h}$ and $A_{R,h,z}$ by $U,P,A$.
For $f\in V_{R,h}$ the exact identity
\begin{equation}\label{eq:compression-telescoping}
 A^nf-PU^nf=\sum_{j=0}^{n-1}A^{n-1-j}PU(P-I)U^jf
\end{equation}
follows by telescoping. All operators on the right are contractions.
For $f=1$, $U^j1=e^{-iz\cdot U_{j,R}}$ under the section law.
The cumulative return tail gives
\[
 \|U^j1-v_{j,L}\|_2
       =\nu_R^*(N_{j,R}>L)^{1/2}\le C e^{(an-cL)/2}.
\]
For $j=0$, $P1=1$ and no error occurs. For $1\le j<n$,
Lemma~\ref{lem:section-grid-reconstruction} applied to
Lemma~\ref{lem:truncated-return-phase-variation} gives
\[
 \|(I-P)U^j1\|_2\le C e^{(an-cL)/2}
       +C\sqrt{h(1+|z|)\exp\{C_1L\log(2+L)\}}.
\]
Take the norm in \eqref{eq:compression-telescoping} and pair with the
unit vector $1$. Since $P1=1$, the pairing of $PU^n1$ is the
unmodified characteristic function. This proves the estimate.
\end{proof}

\paragraph{The remaining power and inversion estimates.}
The operator \eqref{eq:actual-compressed-operator} supplies an explicit
regular reconstruction class, a physical defect comparison, and a
local peripheral resolvent. These are proved for every mesh meeting
\eqref{eq:compressed-peripheral-domain}. Passing to the uncompressed
raw record additionally requires a mesh and time scale satisfying
\eqref{eq:compressed-finite-time-error}, control of the other spectral
arcs, and a quantitative conversion of the resulting contour estimates
into powers. The finite-time consistency statement alone supplies
none of these spectral limits. The full $L^2$ operator is unitary;
its contraction properties must not be confused with the spectral
behavior sought on an anisotropic space.

The physical annulus remains \eqref{eq:peripheral-annulus-scales}:
$2n^{-99/200}\le|z|\le n^{-2/5}$, or
$2n^{1/200}\le|\sqrt n\,z|\le n^{1/10}$ after rescaling.
The compact nonzero and growing roof-frequency regions, the far roof
tail, and the global critical/singular raw density extraction remain
the other terms of Theorem~\ref{thm:LLT}. The finite-count cutoff is
used only in the proved error comparison; the matrix entries and the
characteristic function in that comparison retain the full actual
return record. No conditioning event or denominator is replaced.
